\documentclass[a4paper,12pt]{article}
\usepackage[T2A]{fontenc}
\usepackage[utf8]{inputenc}
\usepackage[russian]{babel}
\usepackage{amsmath,amssymb,mathtools}
\usepackage[scaled=0.96]{PTSans}
\renewcommand{\familydefault}{\sfdefault}
\usepackage{icomma}
\usepackage[a4paper,margin=23mm,top=24mm,bottom=23mm,headheight=24pt]{geometry}
\usepackage{microtype}
\usepackage{tikz}
\usetikzlibrary{arrows.meta,positioning,shapes.geometric,calc}
\usepackage{booktabs,tabularx,array}
\usepackage{enumitem}
\usepackage{fancyhdr}
\usepackage{setspace}
\usepackage{xcolor}
\usepackage{titlesec}
\definecolor{accent}{HTML}{176F70}
\definecolor{accentdark}{HTML}{145253}
\definecolor{textgray}{HTML}{30383B}
\definecolor{muted}{HTML}{67757A}
\definecolor{pale}{HTML}{EAF3F1}
\color{textgray}
\onehalfspacing
\setlength{\parindent}{0pt}
\setlength{\parskip}{5pt}
\setlist[enumerate]{leftmargin=*,itemsep=7pt,topsep=6pt}
\setlist[itemize]{leftmargin=*,itemsep=2pt,topsep=3pt}
\titleformat{\section}{\Large\bfseries\color{accentdark}}{\thesection.}{0.65em}{}
\titlespacing*{\section}{0pt}{15pt}{5pt}
\titleformat{\subsection}{\normalsize\bfseries\color{accent}}{}{0pt}{}
\titlespacing*{\subsection}{0pt}{10pt}{2pt}
\pagestyle{fancy}
\fancyhf{}
\fancyhead[L]{\footnotesize\color{muted}Ксения Клыкова \quad / \quad Математика, ЕГЭ}
\fancyhead[R]{\footnotesize\color{muted}17.10.26}
\fancyfoot[C]{\footnotesize\color{muted}\thepage}
\renewcommand{\headrulewidth}{0.3pt}
\renewcommand{\headrule}{\hbox to\headwidth{\color{pale}\leaders\hrule height \headrulewidth\hfill}}
\newcommand{\key}[1]{\textbf{\color{accentdark}#1}}
\newcommand{\mathnote}[1]{\par\vspace{2pt}{\small\color{muted}#1}\par}
\begin{document}
\begin{titlepage}
\thispagestyle{empty}
\begin{center}
\vspace*{33mm}
{\color{accent}\large МАТЕМАТИКА \,$\cdot$\, ПОДГОТОВКА К ЕГЭ}\par
\vspace{17mm}
{\fontsize{29}{35}\selectfont\bfseries Чек-лист}\par
\vspace{10mm}
{\fontsize{19}{26}\selectfont Повторение задач первой части}\par
\vspace{7mm}
{\large Геометрия, векторы, вероятность,\\ математическое ожидание, степени,\\ движение и текстовые задачи}\par
\vspace{26mm}
{\color{muted}\large \today}\par
\vfill
{\color{accentdark}\normalsize Лёвин Артём Александрович\\[2mm] эксклюзивно для Ксении Клыковой}\par
\vspace{18mm}
\end{center}
\begin{tikzpicture}[remember picture,overlay]
\draw[accent!38,line width=1.2pt]
  ($(current page.south west)+(0,17mm)$)
  .. controls ($(current page.south west)+(55mm,29mm)$)
          and ($(current page.south east)+(-60mm,3mm)$)
  .. ($(current page.south east)+(0,15mm)$);
\end{tikzpicture}
\end{titlepage}

\section{Площадь треугольника}
\key{Проверьте, что Вы помните пять основных формул.}
Для треугольника со сторонами $a,b,c$, площадью $S$, высотой $h_a$, полупериметром $p$, радиусами вписанной и описанной окружностей $r$ и $R$:
\[
\begin{aligned}
 S&=\tfrac12 a h_a, &
 S&=\tfrac12 ab\sin\gamma,\\[2pt]
 S&=\sqrt{p(p-a)(p-b)(p-c)}, &
 p&=\tfrac{a+b+c}{2},\\[2pt]
 S&=pr, &
 S&=\tfrac{abc}{4R}.
\end{aligned}
\]
Угол $\gamma$ в формуле с синусом \emph{заключён между сторонами} $a$ и $b$.

\subsection{Пример}
В равнобедренном треугольнике боковые стороны равны $\sqrt6$, а угол между ними --- $30^\circ$. Тогда
\[
 S=\frac12\cdot\sqrt6\cdot\sqrt6\cdot\sin30^\circ
 =\frac12\cdot6\cdot\frac12=\boxed{\frac32}.
\]
\mathnote{Самопроверка: $\sin30^\circ=\tfrac12$; два множителя $\sqrt6$ дают $6$.}

\section{Векторы: координаты и скалярное произведение}
Если $\vec a=(a_1;a_2)$ и $\vec b=(b_1;b_2)$, то
\[
 k\vec a=(ka_1;ka_2),\qquad
 \vec a\cdot\vec b=a_1b_1+a_2b_2.
\]
Полезное свойство: $(\lambda\vec a)\cdot(\mu\vec b)=\lambda\mu(\vec a\cdot\vec b)$.

\subsection{Пример}
Пусть $\vec a=(-1;4)$, $\vec b=(8;-6)$. Найдём $(2\vec a)\cdot(0{,}5\vec b)$:
\[
 2\vec a=(-2;8),\quad 0{,}5\vec b=(4;-3),\quad
 (-2)\cdot4+8\cdot(-3)=\boxed{-32}.
\]
\mathnote{Не теряйте минусы и заключайте отрицательные координаты в скобки.}

\section{Вероятность: промежуток и два испытания}
\subsection{Вложенные события}
Если $X$ --- целое число пассажиров, то событие «от 21 до 24 включительно» можно представить как разность вложенных событий:
\[
 P(21\le X\le24)=P(X\ge21)-P(X\ge25).
\]
Например, при $P(X>20)=0{,}62$, $P(X>24)=0{,}38$ получаем
\[
 P(21\le X\le24)=0{,}62-0{,}38=\boxed{0{,}24}.
\]
\mathnote{«Больше 24» означает 25 и более; «не менее 24» означает 24 и более.}

\subsection{Два независимых матча}
Команда получает $3$ очка за победу, $1$ за ничью, $0$ за поражение. Пусть вероятности дома: $P(\text{победа})=0{,}6$, $P(\text{ничья})=0{,}2$; в гостях: $0{,}4$ и $0{,}1$. Результаты матчей считаются независимыми.

Для суммы очков \key{не менее четырёх} подходят ровно три пары результатов:
\begin{center}
\renewcommand{\arraystretch}{1.18}
\begin{tabularx}{\linewidth}{@{}X X c c@{}}
\toprule
Дома & В гостях & Очки & Вероятность \\
\midrule
Победа & Победа & $3+3=6$ & $0{,}6\cdot0{,}4=0{,}24$\\
Победа & Ничья & $3+1=4$ & $0{,}6\cdot0{,}1=0{,}06$\\
Ничья & Победа & $1+3=4$ & $0{,}2\cdot0{,}4=0{,}08$\\
\bottomrule
\end{tabularx}
\end{center}
Эти варианты несовместны, поэтому вероятности складываются:
\[
 P(\text{не менее }4\text{ очков})=0{,}24+0{,}06+0{,}08=\boxed{0{,}38}.
\]
\mathnote{Перемножаем вероятности независимых результатов в одной строке; складываем вероятности разных подходящих строк.}

\section{Математическое ожидание}
Для дискретной случайной величины $X$ со значениями $x_i$ и вероятностями $p_i$:
\[
 \boxed{\mathsf E(X)=x_1p_1+x_2p_2+\dots+x_np_n},\qquad
 p_1+\dots+p_n=1.
\]
\subsection{Пример}
Из $100$ покупателей $77$ не купили бутылок, $10$ купили одну, $5$ --- упаковку из четырёх, $8$ --- упаковку из шести:
\begin{center}
\begin{tabular}{@{}lcccc@{}}
\toprule
Число бутылок $x_i$ & $0$ & $1$ & $4$ & $6$\\
Вероятность $p_i$ & $0{,}77$ & $0{,}10$ & $0{,}05$ & $0{,}08$\\
\bottomrule
\end{tabular}
\end{center}
\[
 \mathsf E(X)=0\cdot0{,}77+1\cdot0{,}10+4\cdot0{,}05+6\cdot0{,}08
 =\boxed{0{,}78}.
\]
\mathnote{Математическое ожидание может быть дробным и больше 1. Здесь $0\le\mathsf E(X)\le6$; ограничение $p_i\le1$ относится к вероятностям.}

\section{Степени с одинаковым основанием}
\[
 (a^m)^n=a^{mn},\qquad \sqrt{a^{\,m}}=a^{m/2}\quad(a>0).
\]
Если $a>0$ и $a\ne1$, то из $a^u=a^v$ следует $u=v$.

\subsection{Пример}
Решим $25^{x-1}=\sqrt{125^{x+1}}$. Представим $25=5^2$, $125=5^3$:
\[
 5^{2(x-1)}=5^{\frac32(x+1)}
 \ \Longrightarrow\ 2(x-1)=\frac32(x+1)
 \ \Longrightarrow\ \boxed{x=7}.
\]
\mathnote{Здесь основания степеней сравнивают, а не «делят обе части на 5».}

\section{Движение по квадратичному закону}
Пусть высота тела задаётся функцией $h(t)$ и нужно узнать, \emph{сколько времени} высота была не меньше $H$.
\begin{enumerate}[label=\arabic*. ,itemsep=1pt]
\item Решите неравенство $h(t)\ge H$ при $t\ge0$.
\item Определите моменты пересечения уровня $H$: $t_1\le t_2$.
\item Если движение выше уровня происходит между этими моментами, длительность равна $\Delta t=t_2-t_1$.
\end{enumerate}

\subsection{Пример}
Для $h(t)=-10t^2+26t+1{,}2$ найдём время на высоте не менее $10$ м:
\[
 \begin{aligned}
 -10t^2+26t+1{,}2&\ge10,\\
 10t^2-26t+8{,}8&\le0,\\
 t_1=0{,}4,\quad t_2&=2{,}2.
 \end{aligned}
\]
Так как парабола $10t^2-26t+8{,}8$ направлена ветвями вверх, подходящий промежуток --- между корнями:
\[
 t\in[0{,}4;2{,}2],\qquad \Delta t=2{,}2-0{,}4=\boxed{1{,}8\text{ с}}.
\]
\mathnote{Значение $t_1$ --- первое достижение высоты при подъёме, $t_2$ --- второе при падении.}

\section{Текстовая задача на движение}
\key{Основная связь:} $S=vt$, поэтому $t=\dfrac Sv$.

\subsection{Пример}
Автомобилист и мотоциклист одновременно выехали на расстояние $240$ км. Мотоциклист едет на $12$ км/ч быстрее и прибывает на $40$ минут раньше. Обозначим его скорость через $x$ км/ч.
\begin{center}
\renewcommand{\arraystretch}{1.25}
\begin{tabularx}{\linewidth}{@{}X c c c@{}}
\toprule
Участник & Скорость, км/ч & Время, ч & Путь, км\\
\midrule
Мотоциклист & $x$ & $240/x$ & $240$\\
Автомобилист & $x-12$ & $240/(x-12)$ & $240$\\
\bottomrule
\end{tabularx}
\end{center}
Переводим $40$ минут в часы: $40/60=2/3$. Автомобилист потратил больше времени:
\[
 \frac{240}{x-12}-\frac{240}{x}=\frac23,\qquad x>12.
\]
Умножая на $3x(x-12)$, получаем $x^2-12x-4320=0$. Из двух корней условию $x>12$ удовлетворяет $\boxed{x=72\text{ км/ч}}$.
\mathnote{Проверка: автомобилист ехал $240/60=4$ ч, мотоциклист $240/72=3\tfrac13$ ч. Разница $\tfrac23$ ч.}

\clearpage
\section*{Алгоритм: как составить уравнение в задаче на движение}
\thispagestyle{plain}
\vspace{5mm}
\begin{center}
\tikzset{
 base/.style={align=center,draw=accent!55,text=textgray,rounded corners=3pt,minimum width=115mm,text width=108mm,minimum height=14mm,inner sep=4pt,font=\normalsize},
 terminal/.style={base,draw=accentdark,fill=pale,rounded corners=8mm},
 stepbox/.style={base,fill=white},
 ar/.style={-{Stealth[length=3.2mm]},line width=0.9pt,draw=accent}
}
\begin{tikzpicture}[node distance=10mm]
\node[terminal] (a) {Прочитайте условие, отметьте путь, скорость и время};
\node[stepbox,below=of a] (b) {Обозначьте через $x$ искомую скорость\\и запишите ограничение $x>0$ (при необходимости строже)};
\node[stepbox,below=of b] (c) {Составьте таблицу $v$ --- $t$ --- $S$\\и выразите неизвестное время по формуле $t=S/v$};
\node[stepbox,below=of c] (d) {Переведите минуты в часы и сопоставьте времена\\(кто прибыл раньше, кто потратил больше времени)};
\node[stepbox,below=of d] (e) {Запишите уравнение как разность времён\\и решите его с учётом ОДЗ};
\node[terminal,below=of e] (f) {Проверьте скорости и исходную разницу времён};
\draw[ar] (a)--(b);
\draw[ar] (b)--(c);
\draw[ar] (c)--(d);
\draw[ar] (d)--(e);
\draw[ar] (e)--(f);
\end{tikzpicture}
\end{center}
\clearpage
\section*{Тренировочный блок \quad \color{accent}\large Путь А}
\textcolor{muted}{10 заданий. Решения оформляйте самостоятельно; ответы вынесены на отдельную страницу.}
\begin{enumerate}[label=\textbf{\arabic*.}]
\item В равнобедренном треугольнике боковая сторона равна $6$, а угол при вершине между боковыми сторонами равен $30^\circ$. Найдите площадь.
\item Даны векторы $\vec a=(-3;2)$ и $\vec b=(6;-4)$. Найдите $(2\vec a)\cdot(0{,}5\vec b)$.
\item Вероятность того, что пассажиров в автобусе не менее $17$, равна $0{,}72$, а не менее $23$ --- $0{,}28$. Найдите вероятность того, что пассажиров от $17$ до $22$ включительно.
\item Среди $100$ посетителей магазина $70$ ничего не купили, $15$ купили одну бутылку, $10$ --- упаковку из четырёх, $5$ --- упаковку из шести. Найдите математическое ожидание числа купленных бутылок.
\item Футбольная команда в домашнем матче побеждает с вероятностью $0{,}5$ и играет вничью с вероятностью $0{,}3$; в гостевом --- с вероятностями $0{,}4$ и $0{,}2$. Матчи независимы. За победу дают $3$ очка, за ничью --- $1$. Найдите вероятность набрать в двух матчах не менее $4$ очков.
\item Решите уравнение $25^{x-1}=\sqrt{125^{x+1}}$.
\item Высота мяча задана формулой $h(t)=-5t^2+20t+1$, где $h$ --- в метрах, $t$ --- в секундах. Сколько секунд мяч находится на высоте не менее $16$ м?
\item Высота мяча задана формулой $h(t)=-10t^2+30t+2$. Сколько секунд мяч находится на высоте не менее $22$ м?
\item Мотоциклист и автомобилист одновременно выехали из $A$ в $B$, расстояние $120$ км. Мотоциклист едет на $10$ км/ч быстрее и прибывает на $20$ минут раньше. Обозначив скорость мотоциклиста через $x$, \emph{составьте уравнение}; решать его не требуется.
\item Автомобилист и мотоциклист одновременно выехали из $A$ в $B$, расстояние $240$ км. Мотоциклист едет на $16$ км/ч быстрее и прибывает на $45$ минут раньше. Найдите скорость мотоциклиста.
\end{enumerate}

\clearpage
\section*{Ответы к тренировочному блоку}
\vspace{5mm}
\renewcommand{\arraystretch}{1.50}
\begin{tabularx}{\linewidth}{@{}c X@{}}
\toprule
№ & Ответ\\
\midrule
1 & $9$\\
2 & $-26$\\
3 & $0{,}44$\\
4 & $0{,}85$\\
5 & $0{,}42$\\
6 & $x=7$\\
7 & $2$ с\\
8 & $1$ с\\
9 & $\displaystyle\frac{120}{x-10}-\frac{120}{x}=\frac13,\quad x>10$\\
10 & $80$ км/ч\\
\bottomrule
\end{tabularx}
\vfill
\textcolor{muted}{\small \textbf{Контроль перед завершением:} проверьте единицы измерения, знаки координат, сумму вероятностей и условия допустимости найденных значений.}
\end{document}
